{\small\renewcommand{\arraystretch}{1.25}
  \small
  \renewcommand{\arraystretch}{1.25}
  \begin{longtable}{>{\raggedright\arraybackslash}p{0.2043\textwidth}>{\raggedright\arraybackslash}p{0.7245\textwidth}}
    \toprule
    \textbf{Ký hiệu} & \textbf{Ý nghĩa} \\
    \midrule
    ADC & \textit{Analog-to-Digital Converter} --- bộ chuyển đổi tín hiệu tương tự sang số \\
    GSL & \textit{Gas Source Localisation} --- bài toán định vị nguồn phát khí \\
    I$^{2}$C & \textit{Inter-Integrated Circuit} --- chuẩn giao tiếp hai dây giữa vi điều khiển và cảm biến \\
    IQR & \textit{Interquartile Range} --- khoảng tứ phân vị (dùng trong biểu đồ hộp ở Hình 6.3) \\
    LoRa & \textit{Long Range} --- kỹ thuật điều chế trải phổ cho truyền tin xa, công suất thấp \\
    MOX & \textit{Metal-Oxide Semiconductor} --- cảm biến khí kiểu bán dẫn oxit kim loại \\
    PID & \textit{Photo-Ionisation Detector} --- đầu dò ion hoá quang (loại cảm biến khí chính xác cao) \\
    NDIR & \textit{Non-Dispersive Infrared} --- cảm biến khí hồng ngoại không tán sắc \\
    PWM & \textit{Pulse-Width Modulation} --- điều chế độ rộng xung, dùng để điều khiển tốc độ động cơ \\
    ppm & \textit{parts per million} --- phần triệu, đơn vị nồng độ \\
    SPI & \textit{Serial Peripheral Interface} --- chuẩn giao tiếp nối tiếp đồng bộ bốn dây \\
    SRD & \textit{Short Range Devices} --- băng tần dành cho thiết bị vô tuyến cự ly ngắn \\
    Rs & Điện trở của lớp nhạy khí trong cảm biến MQ-3 tại điều kiện đang đo \\
    R$_{0}$ & Điện trở tham chiếu của cảm biến MQ-3, định nghĩa theo datasheet là điện trở trong 0,4 mg/L cồn \\
    RL & Điện trở tải mắc nối tiếp với Rs trên module MQ-3 \\
    A, B & Hai hệ số của đường đặc tuyến luỹ thừa ppm = A$\cdot$(Rs/R$_{0}$)\textsuperscript{B} \\
    N & Số lần chạy trong một tổ hợp thí nghiệm \\
    \textit{p} & Giá trị \textit{p} của kiểm định thống kê \\
    \bottomrule
  \end{longtable}
}
